\documentclass[../../main.tex]{subfiles}
\begin{document}
% SOURCE p=215; printed=200; object=ch06-interpolation-title; region=top
\section{$L^p$ 空间的插值}\label{section:ch06-lp-interpolation}
% SOURCE p=215; printed=200; object=ch06-interpolation-introduction; region=upper
若 $1\leqslant p<q<r\leqslant\infty$，则 $(L^p\cap L^r)\subset L^q\subset(L^p+L^r)$，因而自然会问：定义在 $L^p+L^r$ 上、在 $L^p$ 与 $L^r$ 上均有界的线性算子 $T$，是否也在 $L^q$ 上有界？答案是肯定的，而且可以作多种推广。本节介绍这方面的两个基本定理：Riesz–Thorin 插值定理与 Marcinkiewicz 插值定理。先讨论前者，其证明依赖下面的复函数论结果。
% SOURCE p=215; printed=200; object=ch06-three-lines; number=6.26; region=middle
\begin{lemma}[三线引理]\label{lemma:ch06-three-lines}
设 $\phi$ 为闭带 $0\leqslant\operatorname{Re}z\leqslant1$ 上的有界连续函数，在带的内部全纯。若在 $\operatorname{Re}z=0$ 上有 $|\phi(z)|\leqslant M_0$，在 $\operatorname{Re}z=1$ 上有 $|\phi(z)|\leqslant M_1$，则对 $0<t<1$，在 $\operatorname{Re}z=t$ 上有
\begin{align*}
|\phi(z)|\leqslant M_0^{1-t}M_1^t.
\end{align*}
\end{lemma}
% SOURCE p=215; printed=200; object=ch06-three-lines-proof; region=middle
\begin{proof}
对 $\varepsilon>0$，令 $\phi_\varepsilon(z)=\phi(z)M_0^{z-1}M_1^{-z}\exp(\varepsilon z(z-1))$。那么 $\phi_\varepsilon$ 满足引理的假设，而两个边界常数都变为 $1$；并且当 $|\operatorname{Im}z|\to\infty$ 时，$|\phi_\varepsilon(z)|\to0$。因此当 $A$ 足够大时，在矩形 $0\leqslant\operatorname{Re}z\leqslant1$、$-A\leqslant\operatorname{Im}z\leqslant A$ 的边界上，$|\phi_\varepsilon(z)|\leqslant1$。由最大模原理，整个带上也有这个界。令 $\varepsilon\to0$，在 $\operatorname{Re}z=t$ 上得到
\begin{align*}
|\phi(z)|M_0^{t-1}M_1^{-t}
=\lim_{\varepsilon\to0}|\phi_\varepsilon(z)|\leqslant1,
\end{align*}
即所求结论。
\end{proof}
% SUPPLEMENT: id=v14-ch06-three-lines-details; parent=ch06-three-lines; kind=proof-detail+motivation+hypotheses; origin=independent-derivation
\begin{remark}
\textbf{边界退化、动机与逐条件说明（重构新增）。}上述复幂先对 $M_0,M_1>0$ 使用实对数定义。若某个 $M_j=0$，先把两个常数都换为 $M_j+\delta$，应用正数情形，再令 $\delta\downarrow0$；因为 $t,1-t>0$，所得界仍有意义，并给出带内的零值结论。

写 $z=x+iy$，则 $\operatorname{Re}(z(z-1))=x(x-1)-y^2$。若 $|\phi|\leqslant B$，在整个闭带上便有
\begin{align*}
|\phi_\varepsilon(x+iy)|
\leqslant B\max(M_0^{-1},M_1^{-1})e^{-\varepsilon y^2}.
\end{align*}
这是关于 $x\in[0,1]$ 的一致衰减，故同一个 $A$ 控制两条横边。在两条竖边上，常数幂恰消掉 $M_j$，而指数因子的模不超过一。任给带内一点，取包含该点的足够大矩形，最大模原理便给出所需界。连续到边界使矩形边界估计可用；内部全纯性使最大模原理适用；整个带上有界使横边可被高斯衰减压下去。不能仅凭两条无限长竖边的界，直接引用有界区域上的最大模原理。

动机是把两侧的不同常数用复幂归一成一，再以小的衰减因子把无界带转化为有限矩形。结论在对数尺度上按横坐标线性插值，这正对应后面 $L^p$ 指数的倒数插值。本证明调用复分析的最大模原理作为先修结果。
\end{remark}
% END SUPPLEMENT: id=v14-ch06-three-lines-details
% SOURCE p=215; printed=200; object=ch06-riesz-thorin; number=6.27; region=lower
\begin{theorem}[Riesz–Thorin 插值定理]\label{theorem:ch06-riesz-thorin}
设 $(X,\mathscr M,\mu)$ 与 $(Y,\mathscr N,\nu)$ 为测度空间，$p_0,p_1,q_0,q_1\in[1,\infty]$。若 $q_0=q_1=\infty$，还假设 $\nu$ 半有限。对 $0<t<1$，定义
\begin{align*}
\frac1{p_t}=\frac{1-t}{p_0}+\frac t{p_1},\qquad
\frac1{q_t}=\frac{1-t}{q_0}+\frac t{q_1}.
\end{align*}
若 $T:L^{p_0}(\mu)+L^{p_1}(\mu)\to L^{q_0}(\nu)+L^{q_1}(\nu)$ 为线性映射，并且
\begin{align*}
\|Tf\|_{q_0}&\leqslant M_0\|f\|_{p_0}\quad(f\in L^{p_0}(\mu)),\\
\|Tf\|_{q_1}&\leqslant M_1\|f\|_{p_1}\quad(f\in L^{p_1}(\mu)),
\end{align*}
则 $\|Tf\|_{q_t}\leqslant M_0^{1-t}M_1^t\|f\|_{p_t}$ 对所有 $f\in L^{p_t}(\mu)$ 成立。
\end{theorem}
% SOURCE p=215,216,217; printed=200-202; object=ch06-riesz-thorin-proof; region=bottom-top
\begin{proof}
先注意，$p_0=p_1=p$ 时，由\cref{proposition:ch06-lp-interpolation}，
\begin{align*}
\|Tf\|_{q_t}\leqslant\|Tf\|_{q_0}^{1-t}\|Tf\|_{q_1}^t
\leqslant M_0^{1-t}M_1^t\|f\|_p.
\end{align*}
因此可以假设 $p_0\ne p_1$，特别地，$p_t<\infty$。记 $\Sigma_X$、$\Sigma_Y$ 为相应空间上在有限测度集之外为零的简单函数空间。由\cref{proposition:ch06-simple-density}，$\Sigma_X$ 包含于所有 $L^p(\mu)$，并在每个有限指数的 $L^p(\mu)$ 中稠密；$Y$ 上也一样。

证明的主要部分是对 $f\in\Sigma_X$ 建立所求估计。由\cref{theorem:ch06-dual-test-criterion}，
\begin{align*}
\|Tf\|_{q_t}=\sup\left\{\left|\int_Y(Tf)g\,\mathrm d\nu\right|:
g\in\Sigma_Y,\ \|g\|_{q_t'}=1\right\},
\end{align*}
其中 $q_t'$ 为 $q_t$ 的共轭指数。注意 $Tf\in L^{q_0}\cap L^{q_1}$，所以除非 $q_0=q_1=\infty$，其非零支集必为 $\sigma$ 有限；在双无穷的情形，假设的半有限性保证对偶测试定理可用。可以设 $f\ne0$，并将它归一为 $\|f\|_{p_t}=1$。于是只须证明：若 $g\in\Sigma_Y$ 且 $\|g\|_{q_t'}=1$，则 $|\int(Tf)g\,\mathrm d\nu|\leqslant M_0^{1-t}M_1^t$。

写 $f=\sum_{j=1}^m c_j\chi_{E_j}$、$g=\sum_{k=1}^n d_k\chi_{F_k}$，其中各 $E_j$、各 $F_k$ 分别互不相交，系数都非零。写成极坐标 $c_j=|c_j|e^{i\theta_j}$、$d_k=|d_k|e^{i\psi_k}$，并令
\begin{align*}
\alpha(z)=(1-z)p_0^{-1}+zp_1^{-1},\qquad
\beta(z)=(1-z)q_0^{-1}+zq_1^{-1}.
\end{align*}
固定 $t\in(0,1)$。因为 $\alpha(t)=p_t^{-1}>0$，可以定义
\begin{align*}
f_z=\sum_{j=1}^m|c_j|^{\alpha(z)/\alpha(t)}e^{i\theta_j}\chi_{E_j}.
\end{align*}
若 $\beta(t)<1$，定义
\begin{align*}
g_z=\sum_{k=1}^n|d_k|^{(1-\beta(z))/(1-\beta(t))}e^{i\psi_k}\chi_{F_k};
\end{align*}
若 $\beta(t)=1$，则对所有 $z$ 取 $g_z=g$。原书接下来假设 $\beta(t)<1$，并把 $\beta(t)=1$ 时的简单修改留给读者。

最后令 $\phi(z)=\int_Y(Tf_z)g_z\,\mathrm d\nu$。于是
\begin{align*}
\phi(z)&=\sum_{j,k}A_{jk}|c_j|^{\alpha(z)/\alpha(t)}
|d_k|^{(1-\beta(z))/(1-\beta(t))},\\
A_{jk}&=e^{i(\theta_j+\psi_k)}\int_Y(T\chi_{E_j})\chi_{F_k}\,\mathrm d\nu.
\end{align*}
因此 $\phi$ 为整函数，在 $0\leqslant\operatorname{Re}z\leqslant1$ 的带上有界。又 $\phi(t)=\int(Tf)g\,\mathrm d\nu$，故由三线引理，只须证明两条边上的界分别为 $M_0,M_1$。对 $s\in\mathbb R$，
\begin{align*}
\alpha(is)&=p_0^{-1}+is(p_1^{-1}-p_0^{-1}),\\
1-\beta(is)&=(1-q_0^{-1})-is(q_1^{-1}-q_0^{-1}),
\end{align*}
所以在相应非零支集上，$|f_{is}|=|f|^{p_t/p_0}$、$|g_{is}|=|g|^{q_t'/q_0'}$。由 Hölder 不等式与归一化，
\begin{align*}
|\phi(is)|\leqslant\|Tf_{is}\|_{q_0}\|g_{is}\|_{q_0'}
\leqslant M_0\|f_{is}\|_{p_0}\|g_{is}\|_{q_0'}=M_0.
\end{align*}
同样的计算给出 $|\phi(1+is)|\leqslant M_1$，从而证明了简单函数上的估计。

由稠密性，$T|_{\Sigma_X}$ 唯一延拓到 $L^{p_t}$，并满足同一个界。还须证明这个延拓就是原来的 $T$。给定 $f\in L^{p_t}$，取 $f_n\in\Sigma_X$，使 $|f_n|\leqslant|f|$ 且 $f_n\to f$ 逐点成立。令
\begin{align*}
E=\{|f|>1\},\quad g=f\chi_E,\quad g_n=f_n\chi_E,
\quad h=f-g,\quad h_n=f_n-g_n.
\end{align*}
重新标记端点后可设 $p_0<p_1$；则 $g\in L^{p_0}$、$h\in L^{p_1}$。原书在此以控制收敛定理得到
\begin{align*}
\|f_n-f\|_{p_t}\to0,\qquad
\|g_n-g\|_{p_0}\to0,\qquad
\|h_n-h\|_{p_1}\to0.
\end{align*}
于是 $Tg_n\to Tg$ 于 $L^{q_0}$，$Th_n\to Th$ 于 $L^{q_1}$。由\cref{exercise:ch06-norm-measure-convergence}，取适当子列后，两者都几乎处处收敛，所以 $Tf_n\to Tf$ 几乎处处。由 Fatou 引理，
\begin{align*}
\|Tf\|_{q_t}&\leqslant\liminf_n\|Tf_n\|_{q_t}
\leqslant\liminf_nM_0^{1-t}M_1^t\|f_n\|_{p_t}\\
&=M_0^{1-t}M_1^t\|f\|_{p_t}.
\end{align*}
证明完毕；无穷端点处所需的逼近细节补充如下。
\end{proof}
% SUPPLEMENT: id=v14-ch06-riesz-thorin-details; parent=ch06-riesz-thorin; kind=omitted-case+proof-detail+motivation+hypotheses; origin=independent-derivation
\begin{remark}
\textbf{端点、逼近与逐条件说明（重构新增）。}首先补全复幂构造。可删去测度为零的简单函数块，使各保留块有正的有限测度。所有幂只作用于正数 $|c_j|,|d_k|$；在支集外定义函数为零，不使用 $0^0$。若 $p_j<\infty$，则
\begin{align*}
\|f_{j+is}\|_{p_j}^{p_j}=\int_X|f|^{p_t}\,\mathrm d\mu=1;
\end{align*}
若 $p_j=\infty$，该函数在支集内模为一，故范数仍为一。对于 $\beta(t)<1$，$g_{j+is}$ 的共轭端点范数完全相同：有限 $q_j'$ 时其 $q_j'$ 次方积分为 $\int|g|^{q_t'}=1$，无穷端点时模为一。因此两条边的估计均已覆盖。$A_{jk}$ 有限，因为 $T\chi_{E_j}\in L^{q_0}$ 且 $\chi_{F_k}\in L^{q_0'}$；有限指数和的模只依赖 $\operatorname{Re}z$，在闭带上一致有界，不需要对参数积分作未验证的微分。

若 $\beta(t)=1$，由 $q_0^{-1},q_1^{-1}\leqslant1$ 与 $0<t<1$ 必有 $q_0=q_1=1$，于是 $q_0'=q_1'=\infty$。取 $g_z=g$ 后，$\|g_z\|_\infty=1$；$\phi(z)$ 仍是有限个整指数函数的线性组合，而两边满足 $|\phi(j+is)|\leqslant M_j$。三线引理照常给出结论。这补全了原书留给读者的分支。若某个 $M_j=0$，三线引理的零常数处理已经适用。若测试单位球面为空，可按非负上确界约定取值零，或改用 $\|g\|_{q_t'}\leqslant1$ 的闭球写公式，避免零空间例外。

\textbf{原延拓证明的无穷输入端点。}当 $p_1=\infty$ 时，模支配的逐点收敛本身不推出 $\|h_n-h\|_\infty\to0$；此处必须进一步选择逼近序列。取简单函数 $s_n$，满足 $|s_n|\leqslant|f|$、$s_n\to f$ 逐点，并在每个模有界区域上一致逼近。这种选择可由\cref{theorem:ch02-simple-approximation}的简单逼近构造得到；例如将模从下方量化、将相位按越来越细的有限分割量化，再把模截于 $n$。令
\begin{align*}
f_n=s_n\chi_{\{|f|>1/n\}}.
\end{align*}
由于 $f\in L^{p_t}$、$p_t<\infty$，该支集的测度不超过 $n^{p_t}\|f\|_{p_t}^{p_t}$，所以 $f_n\in\Sigma_X$。它仍模支配且逐点收敛。在 $E^c=\{|f|\leqslant1\}$ 上，截去的区域误差不超过 $1/n$，其余误差不超过 $\sup_{E^c}|s_n-f|\to0$，从而 $h_n\to h$ 一致。对有限的 $p_0,p_t$，以 $2^{p_0}|g|^{p_0}$、$2^{p_t}|f|^{p_t}$ 作控制函数使用控制收敛；当 $p_1<\infty$ 时，对 $h_n$ 也如此。故原证需要的三个范数收敛均成立。

若输出端点 $q_j<\infty$，范数收敛给出几乎处处收敛的子列；若 $q_j=\infty$，在删去可数个零集后，范数差直接控制逐点差，不必另求子列。先后取子列可使两个端点同时收敛。最后 $q_t<\infty$ 时使用通常的 Fatou 引理；$q_t=\infty$ 时，若 $a=\liminf\|Tf_n\|_\infty<\infty$，先取范数趋于 $a$ 的子列，再在所有本质界同时成立的满测集上取逐点极限，得 $|Tf|\leqslant a$。$a=\infty$ 时不等式自动成立。这证明了无穷范数版本。

动机是把固定的 $f,g$ 放进全纯族，用边界的两个算子估计控制中间配对。倒数指数的线性关系使边界积分恰好变为归一化的 $p_t$、$q_t'$ 次方。$p_j,q_j\geqslant1$ 保证 Hölder 对偶及 Banach 完备性可用；$T$ 在和空间上定义使不同端点的分解给出同一个算子；线性保证 $\phi$ 的全纯性，也保证 $Tf_n=Tg_n+Th_n$。有限测度支集使简单函数同时属于全部端点空间。原定理在双无穷输出时的半有限假设，是此证明通过有限支集测试恢复本质上确界的充分条件；其余情形由 $Tf$ 的 $\sigma$ 有限支集代替，不应擅自把任意两个测度空间都要求成 $\sigma$ 有限。$0<t<1$ 避免端点退化并保证不同输入端点间的 $p_t$ 有限；两个端点本来已在假设中给定。
\end{remark}
% END SUPPLEMENT: id=v14-ch06-riesz-thorin-details
% SOURCE p=217; printed=202; object=ch06-interpolated-norm-log-convex; region=middle
\begin{proposition}[插值算子范数的对数凸性]\label{proposition:ch06-interpolated-norm-log-convex}
在上述定理的条件下，记 $M(t)=\|T\|_{L^{p_t}\to L^{q_t}}$。则 $M(t)\leqslant M_0^{1-t}M_1^t$，而且若 $0<s<t<u<1$、$t=(1-\tau)s+\tau u$，则
\begin{align*}
M(t)\leqslant M(s)^{1-\tau}M(u)^\tau.
\end{align*}
因而 $\log M(t)$ 是 $t$ 的凸函数。
\end{proposition}
% SOURCE p=217; printed=202; object=ch06-interpolated-norm-log-convex-proof; region=middle
\begin{proof}
第一个不等式就是 Riesz–Thorin 定理的结论，其中可以出现严格不等号。把 $s,u$ 对应的两个空间作为新端点，再次应用同一定理，即得第二个不等式，也就是所说的对数凸性。
\end{proof}
% SUPPLEMENT: id=v14-ch06-interpolated-norm-details; parent=ch06-interpolated-norm-log-convex; kind=proof-detail+motivation+hypotheses; origin=independent-derivation
\begin{remark}
\textbf{零算子与条件说明（重构新增）。}新端点的倒数指数仍按 $\tau$ 线性插值，因为 $\alpha,\beta$ 为仿射函数，$\alpha(t)=(1-\tau)\alpha(s)+\tau\alpha(u)$，$\beta$ 也一样。$T$ 已在两个新空间上有界，且定义在它们的和空间上，因此重复使用定理合法。如果新输出端点都是无穷，则原两个输出端点也都是无穷，所需半有限性仍由原假设保证。

严格说，普通实值函数 $\log M$ 要求 $M>0$。若 $p_0\ne p_1$ 且某个内点 $M(t_0)=0$，则 $T$ 在 $\Sigma_X$ 上为零；这些简单函数在每个内部 $L^{p_t}$ 稠密，已证的有界性使 $T$ 在全部这些空间上为零。若 $p_0=p_1$，一个内点范数为零已使 $T$ 在这个共同定义域上为零。因此内部要么全为零，要么全为正；零的情形用始终有意义的乘法不等式，或约定 $\log0=-\infty$。动机是说明所得界可在任意两个中间点继续改进，而不限于最初给定的两个端点常数。
\end{remark}
% END SUPPLEMENT: id=v14-ch06-interpolated-norm-details
% SOURCE p=217; printed=202; object=ch06-sublinear-introduction; region=lower
转向 Marcinkiewicz 定理之前，需要再引入一些术语。设 $T$ 从 $(X,\mathscr M,\mu)$ 上可测函数的某个向量空间 $\mathscr D$，映入 $(Y,\mathscr N,\nu)$ 上全体可测函数的空间。
% SOURCE p=217; printed=202; object=ch06-sublinear-definition; region=lower
\begin{definition}[次线性映射]\label{definition:ch06-sublinear}
若对所有 $f,g\in\mathscr D$ 及 $c>0$，有
\begin{align*}
|T(f+g)|\leqslant|Tf|+|Tg|,\qquad |T(cf)|=c|Tf|,
\end{align*}
则称 $T$ 为\textbf{次线性}的。
\end{definition}
% SUPPLEMENT: id=v14-ch06-sublinear-motivation; parent=ch06-sublinear-definition; kind=motivation+hypotheses; origin=independent-derivation
\begin{remark}
\textbf{动机与逐条件说明（重构新增）。}线性算子满足这些模不等式，但极大算子等非线性对象也满足，所以这一结构扩大了插值的适用范围。向量空间定义域保证和与正倍数仍有定义；输出可测性保证可以谈水平集与弱范数。第一条允许把输入分成大、小两部分后控制输出的水平集；第二条只要求正数倍的模齐次性，不能据此断言 $T(cf)=cTf$ 或任意复数的线性。取 $f=0,c=2$ 可得 $|T0|=0$。在几乎处处等价类的空间使用时，相应等式和不等式均按几乎处处理解。
\end{remark}
% END SUPPLEMENT: id=v14-ch06-sublinear-motivation
% SOURCE p=217; printed=202; object=ch06-strong-type-definition; region=bottom
\begin{definition}[强型]\label{definition:ch06-strong-type}
设 $1\leqslant p,q\leqslant\infty$。若次线性映射 $T$ 的定义域包含 $L^p(\mu)$，它把 $L^p(\mu)$ 映入 $L^q(\nu)$，并且存在常数 $C>0$ 使所有 $f\in L^p(\mu)$ 都满足 $\|Tf\|_q\leqslant C\|f\|_p$，则称 $T$ 为\textbf{强 $(p,q)$ 型}的。
\end{definition}
% SUPPLEMENT: id=v14-ch06-strong-type-motivation; parent=ch06-strong-type-definition; kind=motivation+hypotheses; origin=independent-derivation
\begin{remark}
\textbf{动机与逐条件说明（重构新增）。}第一个指数衡量输入，第二个指数衡量输出；二者所用的测度空间可以不同。定义域包含整个 $L^p$，排除了只对某个测试类写估计而未作延拓的情形。常数必须与 $f$ 无关，才能称为一致有界估计。指数从一开始，使这里使用的是前面构造的范数空间；无穷指数按本质上确界解释。对线性 $T$，这就是通常的有界算子条件；次线性情形仍使用同一不等式，但不称其为线性算子范数。
\end{remark}
% END SUPPLEMENT: id=v14-ch06-strong-type-motivation
% SOURCE p=217; printed=202; object=ch06-weak-type-definition; region=bottom
\begin{definition}[弱型]\label{definition:ch06-weak-type}
设 $1\leqslant p\leqslant\infty$、$1\leqslant q<\infty$。若次线性映射 $T$ 的定义域包含 $L^p(\mu)$，它把 $L^p(\mu)$ 映入弱 $L^q(\nu)$，并且存在常数 $C>0$ 使所有 $f\in L^p(\mu)$ 都满足 $[Tf]_q\leqslant C\|f\|_p$，则称 $T$ 为\textbf{弱 $(p,q)$ 型}的。另外，约定弱 $(p,\infty)$ 型就是强 $(p,\infty)$ 型。
\end{definition}
% SUPPLEMENT: id=v14-ch06-weak-type-motivation; parent=ch06-weak-type-definition; kind=motivation+hypotheses; origin=independent-derivation
\begin{remark}
\textbf{动机与逐条件说明（重构新增）。}有限 $q$ 时，这恰好等价于对每个 $\alpha>0$ 有 $\lambda_{Tf}(\alpha)\leqslant(C\|f\|_p/\alpha)^q$。它只控制超过高度 $\alpha$ 的输出区域，而不直接控制输出的完整 $q$ 次积分。强型由 Chebyshev 不等式蕴含弱型；反向一般不成立。定义仍要求整个输入空间和统一的常数，输出可测使全部水平集有意义。不能把 $q=\infty$ 直接代入有限幂公式，故须单独约定；这一端点约定在下面的截断证明中表现为某个高度以上的水平集为空。
\end{remark}
% END SUPPLEMENT: id=v14-ch06-weak-type-motivation
\subfile{section-05-marcinkiewicz.tex}
\subfile{section-05-kernels.tex}
\subfile{section-05-exercises.tex}
\end{document}
